\chapter{Repository Walkthrough}
\label{chap:repo-walkthrough}
\index{repository walkthrough}\index{hook system}\index{explicit TM API}

This appendix is a guided tour of the companion repository.  Every
\emph{[implementation]} exercise in the book refers to at least one
directory listed here; this chapter maps those references to concrete
paths and explains the role of each component.  Read this appendix before
starting Part~IV or Part~V exercises.

\section{Repository Structure}
\label{sec:repo-structure}
\index{repository!directory layout}

The top-level layout of the repository:

\begin{verbatim}
clang-tm/
+-- backends/tm_impl/        # TM backend implementations (C++)
|   +-- common/              # hook system, region allocator, macros
|   +-- tiny_stm/            # TinySTM WBCTL/WT/WBETL
|   +-- norec/               # NOrec + NOrec-BF
|   +-- tl2/                 # TL2 + TSC-TM
|   +-- romulus/             # version-table OCC
|   +-- leftright/           # global-clock OCC
|   +-- xtm/                 # page-private copy
|   +-- single_global_lock/  # SGL baseline
|   +-- tsx_sgl/             # Intel TSX + SGL fallback
|   +-- spht/                # group-commit persistent HTM
|   +-- calvin/              # deterministic two-phase OCC
|   +-- mvlog/               # multi-version commit-log STM
|   +-- jvstm/               # multi-version OCC (VBox)
|   +-- norec_bf/            # NOrec + Bloom-filter read set
|   +-- tikv/                # distributed TiKV 2PC
|   +-- queue/               # async queue runtime
|   +-- tm_region_allocator/ # mmap region + deferred frees
|   +-- ...                  # 20+ backends total
+-- explicit_api/cpp/        # C++ explicit TM API (TM<T>)
+-- explicit_api/rust/       # Rust workspace (backends + executor)
+-- benchmarks/cpp/          # bank, fuzz-counter, fuzz-bank, STAMP
+-- benchmarks/rust/         # Rust benchmark ports
+-- benchmarks/plugin/       # plugin-instrumented benchmarks
+-- gpu/                     # GPU backends (PR-STM, CSMV, GUST, GPUTX)
+-- tests/                   # test_tx, test_ds, fuzz, plugin tests
+-- plugin/                  # LLVM instrumentation pass + clang-tm driver
+-- docs/proofs/             # TLA+/PlusCal models for all backends
+-- docs/book/               # this book's LaTeX source
+-- simulator/               # discrete-event TM simulator
+-- patches/                 # debug patches, TSX profiling patches
+-- tools/                   # stm_bug_tool, profiling, debug patches
\end{verbatim}

Each directory under \rep{backends/tm\_impl/} contains one or more
backend implementations.  The \emph{common} subdirectory
(\rep{backends/tm\_impl/common/}) is the central hub: it defines the
hook system that connects benchmarks to backends, the region allocator
that provides the TM address space, and the address-classification macros
used by the LLVM pass.  Individual backends do \emph{not} contain
the hook definitions; they only register their implementations at
\texttt{tm\_init()} time.

The \rep{explicit\_api/} tree holds the typed C++ and Rust API layers
that wrap the raw hook function pointers.  The \rep{plugin/} tree
contains the LLVM pass that instruments source code and the
\texttt{clang-tm} driver script.  The \rep{docs/proofs/} directory
holds the TLA+ models verified by TLC (see \S\ref{sec:repo-tla}).

\section{The Hook System}
\label{sec:repo-hooks}
\index{hook system}\index{TMRealHooks}\index{tm\_register\_real\_hooks}

The hook system is the central dispatch mechanism that connects
benchmarks to backends.  It lives entirely in
\rep{backends/tm\_impl/common/}---not inside individual backends.

\subsection{Architecture}

The system has three layers:

\begin{enumerate}
    \item \textbf{Global function-pointer variables} (DATA symbols)
          declared in \rep{backends/tm\_impl/common/tm\_hooks.hpp}.  Each
          TM operation (\texttt{tm\_begin}, \texttt{tm\_end}, seven reads,
          seven writes, \texttt{tm\_malloc}, \texttt{tm\_calloc}, and
          several more) is a \texttt{void (*tm\_op)(...)} variable.
    \item \textbf{Stub defaults} defined in
          \rep{backends/tm\_impl/common/tm\_hooks.cpp}.  At program start,
          every pointer points to a stub that performs plain memory access
          with zero TM overhead.
    \item \textbf{Backend registration}: each backend's
          \texttt{***\_runtime.cpp} populates a \texttt{TMRealHooks}
          struct and calls \texttt{tm\_register\_real\_hooks()} to replace
          the stubs with real implementations.
\end{enumerate}

\subsection{The Registration Pattern}

Every backend follows the same pattern.  The NOrec runtime
(\rep{backends/tm\_impl/norec/NOrec\_runtime.cpp}) is representative:

\begin{lstlisting}[language=C++, caption={Hook registration in the NOrec backend (abridged).}]
// 1. Static implementations of each hook
static void real_tm_begin() {
    if (tm_longjmp_ret == 0) {
        tm_clear_spec_allocs();
        tm_clear_deferred_frees();
        g_in_tx = true;
        norec::begin();
    }
}
// ... real_tm_end, real_tm_read_i4, etc.

// 2. Registration table (designated initializers)
const TMRealHooks g_norec_hooks = {
    .begin    = real_tm_begin,
    .end      = real_tm_end,
    .malloc   = real_tm_malloc,
    .calloc   = real_tm_calloc,
    .read_i4  = real_tm_read_i4,
    .write_i4 = real_tm_write_i4,
    // ... all 24 fields
};

// 3. Called from tm_init()
void tm_init() {
    norec::init();
    tm_register_real_hooks(&g_norec_hooks);
}
\end{lstlisting}

The \texttt{TMRealHooks} struct (\rep{backends/tm\_impl/common/tm\_hooks.hpp:117})
has 24 function-pointer fields covering lifecycle, memory, reads, writes,
and thread-state queries.  Calling \texttt{tm\_register\_real\_hooks}
(\rep{backends/tm\_impl/common/tm\_hooks.cpp:215}) copies the struct
into a module-static \texttt{s\_real\_hooks} variable and calls
\texttt{apply\_hooks\_unlocked()}, which writes each function pointer
into the global DATA variables.

\subsection{Why DATA, Not TEXT}

The LLVM instrumentation pass emits indirect calls through the global
function pointers: \texttt{call *(\%rax)} where \texttt{\%rax} was loaded
from \texttt{tm\_read\_i4}.  Because the pointer is a DATA symbol, the
linker resolves it to the variable's address; the generated code loads
the function address from that variable at call time.  This is why
\texttt{tm\_begin} must be a function \emph{pointer} (DATA), not a
function \emph{declaration} (TEXT).  If it were TEXT, the pass would emit
a direct call to the function's machine code, which on some platforms
collides with the DATA symbol of the same name---see
Appendix~\ref{chap:rust-mistakes} Exercise~3 for the concrete failure.

\subsection{Runtime Backend Swapping}

Because every TM operation dispatches through a global pointer,
\texttt{tm\_swap\_runtime()} (\rep{backends/tm\_impl/common/tm\_hooks.cpp:252})
can replace all hooks atomically under a mutex.  This enables
phase-based TM: run with NOrec, then swap to TinySTM mid-execution, with
all subsequent TM calls going through the new backend.  Benchmarks call
the function pointers---they never call backend functions directly.

\section{The Explicit API}
\label{sec:repo-explicit-api}
\index{explicit TM API}\index{TM<T>}

The C++ explicit API wraps the raw hook function pointers in a typed
class.  The header is
\rep{explicit\_api/cpp/expli\_tm\_api/tm\_api.hpp}.

\subsection{The TM<T> Class}

\begin{lstlisting}[language=C++, caption={\texttt{TM<T>} -- the transactional memory handle.}]
template<typename T>
class TM {
    T *value_ = nullptr;  // allocated via tm_malloc

    T* ptr() {
        if (!value_)
            value_ = (T*)tm_malloc(sizeof(T));
        return value_;
    }
public:
    // ... lifecycle, read, write
    T read() const {
        return tm_type_traits<T>::read(const_cast<T*>(ptr()));
    }
    void write(const T &v) {
        tm_type_traits<T>::write(ptr(), v);
    }
};
\end{lstlisting}

\texttt{TM<T>} is a typed handle.  The value lives in the TM address
space (allocated via \texttt{tm\_malloc}).  The
\texttt{tm\_type\_traits<T>} template dispatches \texttt{read()} and
\texttt{write()} to the correct typed hook---\texttt{tm\_read\_i4} for
\texttt{int32\_t}, \texttt{tm\_read\_i8} for \texttt{int64\_t}, and so
on.  This indirection is the instrumentation point: the LLVM pass
identifies these calls and routes them through the backend's
read-set/write-set tracking.

\subsection{The Retry Loop}

\texttt{TM<T>::transaction(body)} (\rep{explicit\_api/cpp/expli\_tm\_api/tm\_api.hpp:144})
wraps the body in a \texttt{sigsetjmp}/\texttt{siglongjmp} retry loop:

\begin{lstlisting}[language=C++, caption={Safe retry loop in \texttt{TM<T>::transaction}.}]
template<typename F>
static void transaction(F&& body) {
    volatile bool done = false;
    while (!done) {
        sigsetjmp(tm_jmpbuf, 0);
        tm_nested_call_counter = 1;
        tm_set_jmpbuf(&tm_jmpbuf);
        tm_begin();
        body();
        tm_end();
        done = true;
    }
    tm_nested_call_counter = 0;
}
\end{lstlisting}

The \texttt{sigsetjmp} call lives in the same frame as the retry loop, so
the \texttt{jmp\_buf} is always valid when \texttt{tm\_end} triggers an
abort via \texttt{siglongjmp}.  This eliminates the undefined behavior of
older designs where the \texttt{jmp\_buf} resided in a function that had
already returned (\S\ref{sec:tls-retry-buffer-example}).

\section{Benchmark Entry Points}
\label{sec:repo-benchmarks}
\index{benchmarks}\index{bank benchmark}\index{fuzz benchmark}

The bank benchmark (\rep{benchmarks/cpp/bank/bank.cpp}) demonstrates the
explicit API style:

\begin{lstlisting}[language=C++, caption={Bank transfer and worker thread (abridged).}]
void transfer(int src, int dst, int amount) {
    run_tx([&]() {
        int bal = g_bank->accounts[src].balance.read();
        bal -= amount;
        g_bank->accounts[src].balance.write(bal);

        bal = g_bank->accounts[dst].balance.read();
        bal += amount;
        g_bank->accounts[dst].balance.write(bal);
    });
}

void worker_thread(ThreadData &data) {
    expli::TM<int>::thread_init();
    // ... setup ...
    while (!g_stop.load()) {
        transfer(src, dst, 1);
    }
    expli::TM<int>::thread_exit();
}
\end{lstlisting}

The \texttt{transfer} function reads two balances, debits the source, and
credits the destination, all inside one transaction.  The
\texttt{total\_transactional()} function sums every balance inside a TX;
\texttt{total\_non\_transactional()} reads without TM (safe only after all
threads join).  The \texttt{worker\_thread} function shows the per-thread
lifecycle: \texttt{thread\_init} before the loop, \texttt{thread\_exit}
after.

The fuzz benchmarks (\rep{benchmarks/cpp/fuzz\_counter/},
\rep{benchmarks/cpp/fuzz\_bank/}) use the raw hook API with explicit
\texttt{tm\_begin}/\texttt{tm\_end} and a manual retry loop, showing the
alternative style without the \texttt{TM<T>} wrapper.

\section{Queue and Async Execution}
\label{sec:repo-queue}
\index{queue runtime}\index{async execution}

The queue runtime (\rep{backends/tm\_impl/queue/queue\_runtime.h}) provides
asynchronous transaction execution via a thread pool:

\begin{itemize}
    \item \texttt{tm\_enqueue(fn, args)} --- DATA function pointer that
          submits a transaction body for background execution.
    \item \texttt{tm\_wait\_prev\_tx()} --- block until all pending
          transactions submitted by the calling thread complete.
    \item \texttt{tm\_queue\_init(workers, queues)} --- start the thread
          pool with the given worker and queue counts.
\end{itemize}

The Rust workspace provides an equivalent executor at
\rep{explicit\_api/rust/workspace/tm-executor/src/lib.rs}:

\begin{itemize}
    \item \texttt{TxExecutor} trait with \texttt{execute()} and
          \texttt{execute\_async()} methods.
    \item \texttt{InlineExecutor} runs transactions on the calling thread.
    \item \texttt{QueueExecutor} dispatches to a worker pool.
    \item \texttt{spec\_alloc()} / \texttt{spec\_free()} for fast
          per-thread allocation inside TX bodies.
\end{itemize}

\section{TLA+ Model Checking}
\label{sec:repo-tla}
\index{TLA+}\index{model checking}\index{PlusCal}

Every backend has a corresponding TLA+ model in
\rep{docs/proofs/}.  The key models:

\begin{itemize}
    \item \textbf{Volatile backends:} \rep{NOrec.tla}, \rep{SGL.tla},
          \rep{TL2.tla}, \rep{TinySTM\_WBCTL.tla},
          \rep{Leftright.tla}, \rep{Romulus.tla}, \rep{XTM.tla},
          \rep{SwissTM.tla}, \rep{JVSTM.tla}, \rep{MVLog.tla},
          \rep{TSC\_TM.tla}
    \item \textbf{Persistent backends:} \rep{PersistentSGL.tla},
          \rep{DUDETM.tla}, \rep{SPHT.tla}, \rep{NVHTM.tla}
    \item \textbf{Distributed backends:} \rep{DistributedSGL.tla},
          \rep{TiKV.tla}
    \item \textbf{GPU and deterministic:} \rep{GPU\_GUST.tla},
          \rep{GPU\_JVSTM.tla}, \rep{Calvin.tla}
    \item \textbf{Simulator:} \rep{SimEngine.tla}, \rep{TSXSim.tla}
\end{itemize}

Verify a single backend:

\begin{verbatim}
make -C docs/proofs check-one BACKEND=NOrec
\end{verbatim}

Verify all backends and liveness properties:

\begin{verbatim}
make -C docs/proofs check              # safety (all backends)
make -C docs/proofs verify-liveness    # fairness-based liveness
\end{verbatim}

Table~6.1 in Chapter~6 gives the full traceability table linking each
TLA+ model to its C++ implementation, state-space size, and invariant
count.

\section{Debugging and Profiling Tools}
\label{sec:repo-tools}
\index{debugging tools}\index{profiling}

\begin{itemize}
    \item \textbf{Event parser and invariant checker}
          (\rep{tools/stm\_bug\_tool/}): parses TM execution traces and
          checks invariants (money conservation, floor safety) offline.
          Use after a fuzz run finds a failing seed.
    \item \textbf{Debug printf patches}
          (\rep{patches/debug/}): on-demand instrumentation that adds
          \texttt{fprintf} statements to backend internals.  Apply with
          \rep{patches/debug/apply.sh}, revert with
          \rep{patches/debug/revert.sh}.  Useful for tracing read-set
          / write-set sizes, commit decisions, or allocator state during
          active debugging.
    \item \textbf{TSX timing instrumentation}
          (\rep{patches/profile/tsx/}): RDTSC-based profiling of
          \texttt{xbegin}, \texttt{xend}, read/write cycles, and abort
          reason breakdowns.  The \texttt{run\_workflow.sh} script applies
          the patch, builds, runs experiments, and reverts.
    \item \textbf{Discrete-event simulator}
          (\rep{simulator/}): replays instrumented traces through a
          simulated backend with cost-mode timing.  Use \texttt{tm-sim}
          with \texttt{--machine-profile} and \texttt{--clock-mode cost}
          to estimate throughput without hardware.
\end{itemize}

\section{Exercises}
\label{sec:repo-exercises}

\begin{enumerate}
    \item \emph{[implementation]} Build and run the bank benchmark with
          the NOrec backend, then swap to TinySTM at runtime using
          \texttt{tm\_swap\_runtime()}.  Verify money conservation across
          both phases.  The swap test at
          \rep{tests/expli-api/test\_swap\_backends.cpp} demonstrates the
          pattern with two backends linked in a single binary.

    \item \emph{[verification]} Pick any backend from the traceability
          table (Table~6.1), read its PlusCal model in
          \rep{docs/proofs/}, and identify which state variable corresponds
          to which C++ field.  For example, in the NOrec model
          (\rep{NOrec.tla}), the variable \texttt{clock} corresponds to
          the global \texttt{norec::clock} in
          \rep{backends/tm\_impl/norec/NOrec\_globals.hpp}.

    \item \emph{[theory]} Explain why the hook system uses DATA function
          pointers instead of TEXT function calls.  What would break if
          \texttt{tm\_begin} were a regular \texttt{void tm\_begin(void)}
          function?  (Hint: consider the LLVM pass's code generation
          strategy and the linker's symbol resolution, as discussed in
          \S\ref{sec:repo-hooks}.)
\end{enumerate}

\index{repository walkthrough|)}
\index{hook system|)}
\index{explicit TM API|)}
\index{TMRealHooks|)}
\index{tm\_register\_real\_hooks|)}
\index{queue runtime|)}
\index{TLA+|)}
\index{model checking|)}
